\section{Introduction}

\subsection{Motivation}

An industrial robot arm can run the same pick-and-place motion in two scenes that look nearly identical, yet only one is hazardous, for example a vial with its cork missing and a corrosive pictogram, or a knife offered blade-first instead of handle-first. Generic object detection misses these small, local cues, so this thesis builds a scene graph of the work cell from the robot's cameras and an agentic framework that annotates it with evidence-backed risk information: object states, label content, object parts, and the relations that turn them into hazards.

\subsection{Task and Requirements}

This thesis addresses the following tasks and requirements:

\subsubsection{Literature Review}

Review scene graph generation, open-vocabulary perception, LLM- and VLM-based agents, and risk assessment in manipulation, building on the institute's existing risk-annotated scene graph framework.

\subsubsection{Schema Design}

Design a scene graph schema for the pick-and-place work cell (robot, object parts, object states, labels) and a separate schema for risk annotations.

\subsubsection{Scene Graph Generation}

Implement a pipeline that builds and updates the scene graph from the robot's two RGB-D cameras (wrist and third-person) while the pick-and-place task runs, live or from recorded ROS bags.

\subsubsection{Agentic Risk Annotation}

Design and implement an agentic framework that annotates the graph with risk-relevant facts and hazards, supported by evidence and confidence, and exports them as risk triples.

\subsubsection{Evaluation}

Evaluate the framework against baselines on three controlled pick-and-place scenarios, including the chemical vial and cutlery scenarios.

\subsubsection{Documentation}

Document the pipeline, schemas, data, and results so that another researcher can reproduce every result.

\subsection{Contributions}

\subsection{Outline of the Report}
